
\documentclass{article}
\usepackage{colortbl}
\usepackage{makecell}
\usepackage{multirow}
\usepackage{supertabular}

\begin{document}

\newcounter{utterance}

\centering \large Interaction Transcript for game `cladder', experiment `full\_v1.5\_default', episode 8897 with Dialogue Architects\_\_qwen3.5{-}9b\_\_SCoRe\_Qwen3.5{-}9B.
\vspace{24pt}

{ \footnotesize  \setcounter{utterance}{1}
\setlength{\tabcolsep}{0pt}
\begin{supertabular}{c@{$\;$}|p{.15\linewidth}@{}p{.15\linewidth}p{.15\linewidth}p{.15\linewidth}p{.15\linewidth}p{.15\linewidth}}
    \# & \multicolumn{2}{c}{Player} && \multicolumn{2}{c}{Game Master} \\
    \hline

    \theutterance \stepcounter{utterance}  
    & & & \multicolumn{4}{p{0.6\linewidth}}{
        \cellcolor[rgb]{0.9,0.9,0.9}{
            \makecell[{{p{\linewidth}}}]{
                \texttt{\tiny{[P1$\langle$GM]}}
                \texttt{You are an expert in causal inference. The following question is not a typical commonsense query, but rather a meticulously designed question created by a professor specializing in causal inference, intended to assess the students' mastery of the course content.} \\
\texttt{Imagine a self{-}contained, hypothetical world with only the following conditions, and without any unmentioned factors or causal relationships: Yupt has a direct effect on muvq. Zupj has a direct effect on muvq. Method 1: We look directly at how yupt correlates with muvq in general. Method 2: We look at this correlation case by case according to zupj. To understand how yupt affects muvq, is it more correct to use the Method 1 than Method 2?} \\
\texttt{Start your answer with 'Yes' or 'No', followed by additional reasoning or evidence to support your explanation.} \\
            }
        }
    }
     \\ \\

    \theutterance \stepcounter{utterance}  
    & \multicolumn{4}{p{0.6\linewidth}}{
        \cellcolor[rgb]{0.561,0.737,0.561}{
            \makecell[{{p{\linewidth}}}]{
                \texttt{\tiny{[P1$\rangle$GM]}}
                \texttt{No.} \\
\\ 
\texttt{In this hypothetical world, using Method 1 (looking at the general correlation between \$Y\_\{upt\}\$ and \$M\_\{uvq\}\$) is incorrect because it fails to account for **confounding**. Since \$Z\_\{upj\}\$ has a direct effect on \$M\_\{uvq\}\$, \$Z\_\{upj\}\$ acts as a confounder: it influences the outcome (\$M\_\{uvq\}\$) and, by virtue of being a variable in the system alongside \$Y\_\{upt\}\$, creates a non{-}causal association between \$Y\_\{upt\}\$ and \$M\_\{uvq\}\$ if not controlled for.} \\
\\ 
\texttt{Method 1 calculates the marginal correlation \$E{[}M\_\{uvq\} | Y\_\{upt\}{]}\$, which conflates the direct causal effect of \$Y\_\{upt\}\$ with the variation in \$M\_\{uvq\}\$ induced by \$Z\_\{upj\}\$. Unless \$Z\_\{upj\}\$ is independent of \$Y\_\{upt\}\$ (which is not stated and generally assumed false in causal diagrams with common causes or specific structural dependencies), the marginal correlation will be biased.} \\
\\ 
\texttt{Method 2, which examines the correlation "case by case according to \$Z\_\{upj\}\$," describes **stratification** (or conditioning on the confounder). By analyzing the relationship within levels of \$Z\_\{upj\}\$ (i.e., calculating \$E{[}M\_\{uvq\} | Y\_\{upt\}, Z\_\{upj\}=z{]}\$), you block the backdoor path from \$Y\_\{upt\}\$ to \$M\_\{uvq\}\$ through \$Z\_\{upj\}\$. This isolates the direct causal effect of \$Y\_\{upt\}\$ on \$M\_\{uvq\}\$. Therefore, Method 2 is the correct approach to identify the causal effect, while Method 1 yields a biased estimate.} \\
            }
        }
    }
    & & \\ \\

    \theutterance \stepcounter{utterance}  
    & & & \multicolumn{2}{p{0.3\linewidth}}{
        \cellcolor[rgb]{0.9,0.9,0.9}{
            \makecell[{{p{\linewidth}}}]{
                \texttt{\tiny{[GM$|$GM]}}
                \texttt{No} \\
            }
        }
    }
    & & \\ \\

    \theutterance \stepcounter{utterance}  
    & & & \multicolumn{2}{p{0.3\linewidth}}{
        \cellcolor[rgb]{0.9,0.9,0.9}{
            \makecell[{{p{\linewidth}}}]{
                \texttt{\tiny{[GM$|$GM]}}
                \texttt{yes} \\
            }
        }
    }
    & & \\ \\

    \theutterance \stepcounter{utterance}  
    & & & \multicolumn{2}{p{0.3\linewidth}}{
        \cellcolor[rgb]{0.9,0.9,0.9}{
            \makecell[{{p{\linewidth}}}]{
                \texttt{\tiny{[GM$|$GM]}}
                \texttt{game\_result = LOSE} \\
            }
        }
    }
    & & \\ \\

\end{supertabular}
}

\end{document}
